\chapter{Valider le modèle}

\section{Deux niveaux de validation}
Un modèle construit sur des mesures les reproduit toujours bien : ce n'est pas
une preuve. Deux contrôles indépendants sont donc menés.
\begin{enumerate}
\item La \textbf{validation croisée} vérifie que le modèle prédit des images
qu'il n'a pas vues. Les 49 scènes sont réparties en deux groupes (une sur
deux) ; on construit les tables sur un groupe et on mesure l'erreur sur
l'autre, puis on inverse.
\item La \textbf{campagne de rendu} mesure le résultat réel d'\IPR{}, de la
lecture du RAW à l'image finale, contre les exports Lightroom. Elle contrôle
tout ce que la validation croisée ignore : conversions, rendu de base,
compression de sortie.
\end{enumerate}

\section{L'indicateur}
Pour chaque scène, on calcule la variation moyenne de chroma
$\overline{C_{\text{variante}}/C_{\text{neutre}}}-1$ sur deux groupes de pixels :
les couleurs peu saturées ($s<0{,}3$) et les couleurs saturées ($s>0{,}5$).
L'écart est la différence entre le modèle (ou \IPR{}) et Lightroom ; on en
retient la médiane sur les scènes, en valeur absolue, et le 90\textsuperscript{e}
centile, qui renseigne sur les scènes difficiles.

\section{Validation croisée}
Écart médian de variation de chroma (peu saturées / saturées) sur les scènes
exclues de la construction :

\begin{center}
\begin{tabular}{lrrr}
\toprule
Valeur & Ancien modèle A1/A2 & Tables & 90\textsuperscript{e} centile des tables\\
\midrule
$-100$ & 0,020 / 0,122 & 0,002 / 0,006 & 0,007 / 0,026\\
$-50$ & 0,047 / 0,030 & 0,003 / 0,005 & 0,007 / 0,020\\
$-25$ & 0,044 / 0,025 & 0,002 / 0,004 & 0,005 / 0,013\\
$+25$ & 0,079 / 0,076 & 0,004 / 0,005 & 0,012 / 0,015\\
$+50$ & 0,088 / 0,127 & 0,012 / 0,012 & 0,030 / 0,035\\
$+100$ & 0,407 / 0,263 & 0,072 / 0,035 & 0,153 / 0,087\\
\bottomrule
\end{tabular}
\end{center}

Les tables divisent l'erreur par 5 à 20 selon la valeur. Le point faible
reste $+100$ sur les couleurs peu saturées, où la réponse de Lightroom est la
plus forte et la plus variable.

\section{Campagne de rendu}
Variation de chroma médiane sur les scènes, Lightroom puis \IPR{}, et écart
médian par scène (campagne \texttt{mesure-couleur-v008}, 637 états mesurés,
aucune exclusion) :

\begin{center}\small
\begin{tabular}{lrrrr}
\toprule
Valeur & Peu saturées & Saturées & Clarté $Y$ & Écart\\
 & LR / \IPR{} & LR / \IPR{} & LR / \IPR{} & peu / saturées\\
\midrule
$-100$ & $-0{,}97$ / $-0{,}97$ & $-0{,}79$ / $-0{,}79$ & 0,88 / 0,87 & 0,004 / 0,011\\
$-50$ & $-0{,}53$ / $-0{,}53$ & $-0{,}40$ / $-0{,}39$ & 0,95 / 0,95 & 0,003 / 0,004\\
$-25$ & $-0{,}28$ / $-0{,}28$ & $-0{,}19$ / $-0{,}19$ & 0,98 / 0,98 & 0,003 / 0,003\\
$+25$ & $+0{,}11$ / $+0{,}11$ & $+0{,}05$ / $+0{,}06$ & 1,00 / 1,00 & 0,014 / 0,003\\
$+50$ & $+0{,}26$ / $+0{,}28$ & $+0{,}11$ / $+0{,}12$ & 1,00 / 1,00 & 0,031 / 0,007\\
$+100$ & $+0{,}90$ / $+0{,}91$ & $+0{,}22$ / $+0{,}24$ & 0,98 / 0,98 & 0,143 / 0,016\\
\bottomrule
\end{tabular}
\end{center}

Avec l'ancien modèle, la même campagne donnait des écarts de 0,31 / 0,09 à
$+100$, 0,14 / 0,06 à $+50$, et 0,01 / 0,13 à $-100$, sans baisse de clarté au
négatif.

\section{Le cas des grands ciels bleus}
Mesuré par famille de teinte, le gain de chroma d'\IPR{} à $+100$ restait plus
faible que celui de Lightroom sur les bleus (1,66 contre 2,78 pour le
bleu-cyan). Trois vérifications ont localisé la cause.

\begin{enumerate}
\item \textbf{Le modèle est bien appliqué.} Sur les pixels bleus d'\IPR{}, la
table prédit un gain de 1,667 et le rendu donne 1,670 : la compression de
sortie ne retire presque rien (elle concerne 25\,\% de ces pixels, et y réduit
le gain de 2,06 à 1,94).
\item \textbf{Le point de départ est identique.} Les bleus neutres d'\IPR{} et de
Lightroom ont la même chroma (0,055 contre 0,056 en médiane), la même
saturation et la même clarté, scène par scène.
\item \textbf{L'écart vient d'une scène.} La chroma finale des bleus à $+100$
est de 0,118 contre 0,124 en médiane (5\,\% de moins), mais la scène 6V1A9141,
un grand ciel bleu saturé qui couvre 82\,\% de l'image, atteint 0,164 contre
0,196 ($-16\,\%$). Elle pèse lourd dans les statistiques par famille.
\end{enumerate}

Sur ce type de ciel, Lightroom pousse plus fort que son comportement médian.
C'est la trace d'une dépendance à l'image que les 2\,\% de variance « scène »
du chapitre 3 laissaient entrevoir. La reproduire demanderait une variable
supplémentaire, par exemple l'étendue de la zone de couleur, ou un traitement
local ; ce n'est pas fait.

\resultat{Un écart global peut cacher un cas particulier très présent dans
l'échantillon. Avant de corriger un modèle, il faut vérifier qu'il est bien
appliqué, que le point de départ est le même, puis isoler les scènes qui
portent l'écart.}

\section{Historique des versions du modèle}
\begin{center}
\begin{tabular}{lp{11cm}}
\toprule
Version & Changement\\
\midrule
v001 & tables brutes de gain de chroma et de clarté\\
v002 & monotonie imposée par régression isotone\\
v003 & facteur quasi-gris, un par sens du curseur\\
v004 & facteur quasi-gris par valeur (corrige $-25$ et $-50$)\\
v005 & table de rotation de teinte\\
v006 & repli de la rotation qui conserve la teinte (version retenue)\\
\bottomrule
\end{tabular}
\end{center}

Chaque version est une preuve JSON horodatée et immuable
(\texttt{modele-vibrance-lr-vNNN}), qui renvoie à celle qu'elle remplace. Les
campagnes de rendu \texttt{mesure-couleur-v004} à \texttt{v008} mesurent
respectivement l'ancien modèle puis les versions v002, v003, v004 et v006.
